\subsection{REGULAR ELEMENTS}

By Cor. 4 of Th. 2 of §2, no. 2, the regular elements \(g\) of \(G\) can be characterized by either of the following properties:

\begin{enumerate}
\def\labelenumi{\alph{enumi})}
\item
  The subalgebra of g fixed by Ad g is a Cartan subalgebra.
\item
  \(Z(g)_{0}\) is a maximal torus of G.
\end{enumerate}

The set of regular elements of \(\mathrm{G}\) is open and dense in \(\mathrm{G}\) .

In the remainder of this paragraph, we denote by \(G_{r}\) (resp. \(T_{r}\) ) the set of points of G (resp. T) that are regular in G. An element g of G belongs to \(T_{r}\) if and only if \(Z(g)_{0}\) is equal to T; every element of \(G_{r}\) is conjugate to an element of \(T_{r}\) (§2, no. 2, Th. 2).

An element t of T belongs to \(T_{r}\) if and only if \(t^{\alpha} \neq 1\) for every root \(\alpha \in \mathrm{R}(\mathrm{G}, \mathrm{T})\) ; consequently, \(T - T_{r}\) is the union of the subtori Ker \(\alpha\) for \(\alpha\) in \(\mathrm{R}(\mathrm{G}, \mathrm{T})\) .

PROPOSITION 1. Put \(n = \dim G\) . There exists a compact real-analytic manifold \(V\) of dimension \(n - 3\) and an analytic map \(\varphi : V \to G\) whose image is \(G - G_r\) .

Let \(\alpha\in\mathrm{R}(\mathrm{G},\mathrm{T})\) ; put \(\mathrm{V}_{\alpha}=(\mathrm{G}/\mathrm{Z}(\mathrm{Ker}\alpha))\times(\mathrm{Ker}\alpha)\) , and let \(\varphi_{\alpha}\) be the morphism from \(V_{\alpha}\) to G such that, for all \(g\in G\) and all \(t\in\operatorname{Ker}\alpha\) , we have \(\varphi_{\alpha}(\bar{g},t) = gtg^{-1}\) (we denote by \(\bar{g}\) the coset of \(g\) modulo \(\mathrm{Z}(\mathrm{Ker}\alpha)\) ). Then \(\mathrm{V}_{\alpha}\) is a compact real-analytic manifold of dimension

\[
\begin{array}{c} \dim V _ {\alpha} = \dim G - \dim Z (\operatorname{Ker} \alpha) + \dim \operatorname{Ker} \alpha \\ = n - (\dim T + 2) + (\dim T - 1) = n - 3 \end{array}
\]

(§4, no. 5, Th. 1); \(\varphi_{\alpha}\) is a morphism of real-analytic manifolds, and the image of \(\varphi_{\alpha}\) consists of the elements of G conjugate to an element of \(\mathrm{Ker} \alpha\) . It now suffices to take V to be the sum of the manifolds \(V_{\alpha}\) , and \(\varphi\) to be the morphism inducing \(\varphi_{\alpha}\) on each \(V_{\alpha}\) .

Remark. Call an element g of G very regular if \(Z(g)\) is a maximal torus of G. If \(g \in T\) , g is very regular if and only if \(w(g) \neq g\) for every non-identity element w of \(W_{G}(T)\) (§4, no. 7, Prop. 14). Thus, the set of very regular elements is a dense open subset of G (§2, no. 5, Cor. 2 of Prop. 5).

